\section{Related Work}
\label{sec:related_work}

% -----------------------------------------------------------------------
\subsection{Graph-Based Fact Verification}
\label{sec:related_work:graph}

Early graph-based approaches to fact verification represented evidence as
graphs over retrieved sentences or entities and used graph attention
mechanisms to aggregate evidence.  GEAR~\cite{zhou2019gear} constructs an
evidence graph in which each retrieved sentence is a node, with edges
capturing textual similarity, and uses graph attention propagation to
produce a claim verdict.  KGAT~\cite{liu2020kgat} extends this with a
knowledge-graph attention network that explicitly models fine-grained
semantic relationships between claim tokens and evidence tokens via an
interaction graph.  Both systems improve substantially over sentence-level
baselines on FEVER~\cite{thorne2018fever}, but they operate on a fixed,
retrieved evidence set: neither performs structured multi-hop traversal, and
neither expresses calibrated uncertainty over the verdict.

More recent work has moved toward structured retrieval with graph
representations.  GraphCheck~\cite{guo2024graphcheck} uses entity and
relation graphs built from the generated text to detect hallucinations, but
focuses on generation faithfulness rather than claim verification over an
external corpus.  STRIVE~\cite{pan2023strive} introduces structured
retrieval with iterative evidence refinement for multi-step verification.
AFEV~\cite{hu2023afev} applies adversarial training to improve robustness
of fact-checking models on challenging claim distributions.  Cognition
differs from all of these in that its graph is constructed dynamically from
the cassette store per query, requires no pre-defined entity ontology, and
represents evidence as typed infons with explicit polarity and alignment
scores rather than as plain retrieved passages.

% -----------------------------------------------------------------------
\subsection{Selective Generation and Abstention}
\label{sec:related_work:selective}

Calibrated abstention — declining to answer when evidence is absent — has
received increasing attention in both generation and classification
settings.  \citet{joren2025sufficient} study the selective generation
problem: given a context and a question, when should a model generate
versus abstain?  They show that standard LLMs do not reliably abstain when
the context is insufficient.  SelectLLM~\cite{ko2024selectllm} proposes a
learned selection mechanism that routes queries to different-sized LLMs
based on estimated difficulty, with abstention as an option.  Both lines
of work treat abstention as an output-level decision applied post-hoc to a
neural model.  Cognition instead grounds abstention in the DS vacuous mass
$m(\Theta)$: the decision to abstain is structurally enforced by the
combination rule when evidence is absent, rather than calibrated after the
fact.

% -----------------------------------------------------------------------
\subsection{Calibration and Uncertainty Estimation}
\label{sec:related_work:calibration}

\citet{guo2017calibration} demonstrated that modern neural classifiers are
systematically overconfident and that temperature scaling provides a simple
post-hoc calibration.  Evidential Deep Learning~\cite{sensoy2018evidential}
places a Dirichlet prior over class probabilities to produce uncertainty
estimates that are theoretically grounded but require end-to-end training.
R-GCN~\cite{schlichtkrull2018rgcn} extends graph convolutional networks to
relational data, providing the representational backbone for many downstream
fact-checking and knowledge-graph completion systems.  Cognition's DS mass
is not trained for calibration; it emerges directly from the evidence
combination rule, avoiding the distribution-shift fragility of post-hoc
calibration layers.

% -----------------------------------------------------------------------
\subsection{Benchmark Datasets}
\label{sec:related_work:datasets}

FEVER~\cite{thorne2018fever} established the standard three-class
(SUPPORTS / REFUTES / NEI) formulation and remains the most widely used
benchmark for fact verification.  AVeriTeC~\cite{schlichtkrull2023averitec}
extends this to real-world web evidence, requiring systems to retrieve
evidence from live URLs rather than from a pre-indexed Wikipedia dump.  The
AVeriTeC shared task~\cite{averitec2024sharedtask} introduced standardised
evaluation scripts and a competitive leaderboard.  SciFact~\cite{wadden2020scifact}
restricts claims to the biomedical literature with rationale-level
annotation, stressing domain-specific retrieval.  We evaluate on all three
to ensure results are not specific to a single corpus topology.

% -----------------------------------------------------------------------
\subsection{What Cognition Provides That Prior Work Does Not}
\label{sec:related_work:differentiators}

No prior system combines all three of the following properties.
\textbf{Schema-free operation}: Cognition requires no ontology
pre-specification.  Claims are grounded via SPLADE anchor encoding;
connective predicates are autodiscovered from the data.  Graph-based
systems such as GEAR and KGAT presuppose a fixed evidence graph structure
or a pre-indexed knowledge graph.  \textbf{Immutable cassette storage with
time-travel}: Cognition's cassette substrate stores every infon as a
content-addressed, append-only record.  Old snapshots remain queryable;
the audit trail is cryptographically verifiable.  No prior fact-checking
system provides this property.  \textbf{DS mass with explicit $m(\Theta)$
as a first-class output}: calibrated abstention in Cognition is not a
post-hoc calibration layer applied to a neural softmax — it is native to
the inference rule.  When the corpus provides no evidence, $m(\Theta) \to
1.0$ by construction.  Evidential Deep Learning~\cite{sensoy2018evidential}
shares the goal of principled uncertainty but requires full end-to-end
training on labelled data with uncertainty targets; Cognition's uncertainty
requires only the evidence combination rule.
